% GENERATED FILE. Edit canonical lessons, then run npm run generate:books.
% canonical-source-hash: fnv1a64:0d81991942b616b8
% canonical-lessons: SA-C530-bhojyam, SA-C530-peyam, SA-C530-laddukah, SA-C530-panakam, SA-C530-paryusitah

\chapter{Planning a Feast}
\label{ch:sa-planning-a-feast}
\hlchaptermodality{530}

\begin{chapteropening}
\textbf{By the end of this chapter:} \emph{I can plan a feast and say what we will serve if guests come.}

Complete the last lesson of chapter 530: I can plan a feast and say what we will serve if guests come.
\end{chapteropening}

\section[\texorpdfstring{\sk{भोज्यम्} (bhojyam)}{भोज्यम् (bhojyam)}]{\sk{भोज्यम्} (bhojyam) — eatables; food}

\label{lesson:SA-C530-bhojyam}



\begin{quote}
Before the new one: say the Sanskrit for hurries along, then the Sanskrit for haste.
\end{quote}

\subsection*{You'll want to know: \sk{भोज्यम्}}
\textbf{\sk{भोज्यम्}} — \emph{bhojyam} — \textquotedblleft{}eatables; food\textquotedblright{}.

\begin{quote}
\sk{यदि} \sk{अतिथयः} \sk{आगच्छन्ति}, \sk{तर्हि} \sk{बहु} \sk{भोज्यं} \sk{पचामः।}
\end{quote}

\begin{quote}
\emph{yadi atithayaḥ āgacchanti, tarhi bahu bhojyaṃ pacāmaḥ.}
\end{quote}

\begin{quote}
If the guests come, we will cook a lot of food.
\end{quote}

A neuter noun.

\begin{cousinweb}
\textbf{\sk{भोज्यम्}} (\emph{bhojyam}) comes from the root √bhuj, \textquotedblleft{}to enjoy, to eat\textquotedblright{}, with the ending \emph{-ya}; MW gives \emph{bhojya} as eatable, what is eaten (MW).
\end{cousinweb}

\begin{grammarlens}[title={What you've built}]
One more word for saying what happens if.
\end{grammarlens}

\subsection*{Your turn}
\begin{itemize}
\raggedright
  \item \emph{Say it:} \emph{bhojyam}
  \item \emph{Say it:} \emph{bhojyam}, once more
  \item \emph{Say it:} \emph{yadi atithayaḥ āgacchanti, tarhi bahu bhojyaṃ pacāmaḥ.} — If the guests come, we will cook a lot of food.
  \item \emph{Recall:} say the Sanskrit for hurries along, then the Sanskrit for haste, then say \emph{bhojyam} again
\end{itemize}

\begin{tcolorbox}[breakable,colback=teal!4,colframe=teal!35!black,title={Before you move on}]
What does \sk{भोज्यम्} mean? (\textquotedblleft{}Eatables; food\textquotedblright{}.) Say it once more.
\end{tcolorbox}
\section[\texorpdfstring{\sk{पेयम्} (peyam)}{पेयम् (peyam)}]{\sk{पेयम्} (peyam) — a drink}

\label{lesson:SA-C530-peyam}



\begin{quote}
Before the new one: say the Sanskrit for haste, then the Sanskrit for eatables.
\end{quote}

\subsection*{You'll want to know: \sk{पेयम्}}
\textbf{\sk{पेयम्}} — \emph{peyam} — \textquotedblleft{}a drink\textquotedblright{}.

\begin{quote}
\sk{यदि} \sk{ग्रीष्मः} \sk{अस्ति}, \sk{तर्हि} \sk{शीतलं} \sk{पेयं} \sk{ददातु।}
\end{quote}

\begin{quote}
\emph{yadi grīṣmaḥ asti, tarhi śītalaṃ peyaṃ dadātu.}
\end{quote}

\begin{quote}
If it is summer, please give a cold drink.
\end{quote}

A neuter noun.

\begin{cousinweb}
\textbf{\sk{पेयम्}} (\emph{peyam}) comes from the root √pā, \textquotedblleft{}to drink\textquotedblright{}, which also gives \emph{pibati}; MW gives \emph{peya} as to be drunk, drinkable (MW).
\end{cousinweb}

\begin{grammarlens}[title={What you've built}]
One more word for saying what happens if.
\end{grammarlens}

\subsection*{Your turn}
\begin{itemize}
\raggedright
  \item \emph{Say it:} \emph{peyam}
  \item \emph{Say it:} \emph{peyam}, once more
  \item \emph{Say it:} \emph{yadi grīṣmaḥ asti, tarhi śītalaṃ peyaṃ dadātu.} — If it is summer, please give a cold drink.
  \item \emph{Recall:} say the Sanskrit for haste, then the Sanskrit for eatables, then say \emph{peyam} again
\end{itemize}

\begin{tcolorbox}[breakable,colback=teal!4,colframe=teal!35!black,title={Before you move on}]
What does \sk{पेयम्} mean? (\textquotedblleft{}A drink\textquotedblright{}.) Say it once more.
\end{tcolorbox}
\section[\texorpdfstring{\sk{लड्डुकः} (laḍḍukaḥ)}{लड्डुकः (laḍḍukaḥ)}]{\sk{लड्डुकः} (laḍḍukaḥ) — a laddu}

\label{lesson:SA-C530-laddukah}



\begin{quote}
Before the new one: say the Sanskrit for eatables, then the Sanskrit for a drink.
\end{quote}

\subsection*{You'll want to know: \sk{लड्डुकः}}
\textbf{\sk{लड्डुकः}} — \emph{laḍḍukaḥ} — \textquotedblleft{}a laddu\textquotedblright{}.

\begin{quote}
\sk{यदि} \sk{उत्सवः} \sk{भवति}, \sk{तर्हि} \sk{माता} \sk{लड्डुकान्} \sk{पचति।}
\end{quote}

\begin{quote}
\emph{yadi utsavaḥ bhavati, tarhi mātā laḍḍukān pacati.}
\end{quote}

\begin{quote}
If there is a festival, my mother makes laddus.
\end{quote}

A masculine noun.

\begin{cousinweb}
\textbf{\sk{लड्डुकः}} (\emph{laḍḍukaḥ}) is the word MW gives, beside \emph{laḍḍu}, for a sweet made of ground pulse or flour with sugar and spice, fried in ghee; MW gives no derivation, so its origin is uncertain (MW).
\end{cousinweb}

\begin{grammarlens}[title={What you've built}]
One more word for saying what happens if.
\end{grammarlens}

\subsection*{Your turn}
\begin{itemize}
\raggedright
  \item \emph{Say it:} \emph{laḍḍukaḥ}
  \item \emph{Say it:} \emph{laḍḍukaḥ}, once more
  \item \emph{Say it:} \emph{yadi utsavaḥ bhavati, tarhi mātā laḍḍukān pacati.} — If there is a festival, my mother makes laddus.
  \item \emph{Recall:} say the Sanskrit for eatables, then the Sanskrit for a drink, then say \emph{laḍḍukaḥ} again
\end{itemize}

\begin{tcolorbox}[breakable,colback=teal!4,colframe=teal!35!black,title={Before you move on}]
What does \sk{लड्डुकः} mean? (\textquotedblleft{}A laddu\textquotedblright{}.) Say it once more.
\end{tcolorbox}
\section[\texorpdfstring{\sk{पानकम्} (pānakam)}{पानकम् (pānakam)}]{\sk{पानकम्} (pānakam) — a sherbet; a cool drink}

\label{lesson:SA-C530-panakam}



\begin{quote}
Before the new one: say the Sanskrit for a drink, then the Sanskrit for a laddu.
\end{quote}

\subsection*{You'll want to know: \sk{पानकम्}}
\textbf{\sk{पानकम्}} — \emph{pānakam} — \textquotedblleft{}a sherbet; a cool drink\textquotedblright{}.

\begin{quote}
\sk{यदि} \sk{अतिथिः} \sk{श्रान्तः} \sk{अस्ति}, \sk{तर्हि} \sk{तस्मै} \sk{पानकं} \sk{ददातु।}
\end{quote}

\begin{quote}
\emph{yadi atithiḥ śrāntaḥ asti, tarhi tasmai pānakaṃ dadātu.}
\end{quote}

\begin{quote}
If the guest is tired, please give him a sherbet.
\end{quote}

A neuter noun.

\begin{cousinweb}
\textbf{\sk{पानकम्}} (\emph{pānakam}) is built on \emph{pāna}, \textquotedblleft{}drinking\textquotedblright{}, from the root √pā, \textquotedblleft{}to drink\textquotedblright{}; MW gives \emph{pānaka} as a draught, a drink, a beverage (MW).
\end{cousinweb}

\begin{grammarlens}[title={What you've built}]
One more word for saying what happens if.
\end{grammarlens}

\subsection*{Your turn}
\begin{itemize}
\raggedright
  \item \emph{Say it:} \emph{pānakam}
  \item \emph{Say it:} \emph{pānakam}, once more
  \item \emph{Say it:} \emph{yadi atithiḥ śrāntaḥ asti, tarhi tasmai pānakaṃ dadātu.} — If the guest is tired, please give him a sherbet.
  \item \emph{Recall:} say the Sanskrit for a drink, then the Sanskrit for a laddu, then say \emph{pānakam} again
\end{itemize}

\begin{tcolorbox}[breakable,colback=teal!4,colframe=teal!35!black,title={Before you move on}]
What does \sk{पानकम्} mean? (\textquotedblleft{}A sherbet; a cool drink\textquotedblright{}.) Say it once more.
\end{tcolorbox}
\section[\texorpdfstring{\sk{पर्युषितः} (paryuṣitaḥ)}{पर्युषितः (paryuṣitaḥ)}]{\sk{पर्युषितः} (paryuṣitaḥ) — stale}

\label{lesson:SA-C530-paryusitah}



\begin{quote}
Before the new one: say the Sanskrit for a laddu, then the Sanskrit for a sherbet.
\end{quote}

\subsection*{You'll want to know: \sk{पर्युषितः}}
\textbf{\sk{पर्युषितः}} — \emph{paryuṣitaḥ} — \textquotedblleft{}stale\textquotedblright{}.

\begin{quote}
\sk{यदि} \sk{अन्नं} \sk{पर्युषितम्} \sk{अस्ति}, \sk{तर्हि} \sk{अहं} \sk{न} \sk{खादामि।}
\end{quote}

\begin{quote}
\emph{yadi annaṃ paryuṣitam asti, tarhi ahaṃ na khādāmi.}
\end{quote}

\begin{quote}
If the food is stale, I do not eat it.
\end{quote}

\begin{cousinweb}
\textbf{\sk{पर्युषितः}} (\emph{paryuṣitaḥ}) is built on \emph{pari} and the root √vas, \textquotedblleft{}to stay, to pass the night\textquotedblright{}, as its past participle: kept overnight; MW gives \emph{paryuṣita} as not fresh, stale (MW).
\end{cousinweb}

\begin{grammarlens}[title={What you've built}]
One more word for saying what happens if.
\end{grammarlens}

\subsection*{Your turn}
\begin{itemize}
\raggedright
  \item \emph{Say it:} \emph{paryuṣitaḥ}
  \item \emph{Say it:} \emph{paryuṣitaḥ}, once more
  \item \emph{Say it:} \emph{yadi annaṃ paryuṣitam asti, tarhi ahaṃ na khādāmi.} — If the food is stale, I do not eat it.
  \item \emph{Recall:} say the Sanskrit for a laddu, then the Sanskrit for a sherbet, then say \emph{paryuṣitaḥ} again
\end{itemize}

\begin{tcolorbox}[breakable,colback=teal!4,colframe=teal!35!black,title={Before you move on}]
What does \sk{पर्युषितः} mean? (\textquotedblleft{}Stale\textquotedblright{}.) Say it once more.
\end{tcolorbox}
